% Расчётные формулы приложения survivability
% Марковская модель «отказ–восстановление» оценки живучести узла.
\documentclass[12pt]{article}
\usepackage[T2A]{fontenc}
\usepackage[utf8]{inputenc}
\usepackage[russian]{babel}
\usepackage{amsmath, amssymb}
\usepackage[a4paper, margin=2.5cm]{geometry}

\title{Расчётные формулы оценки живучести узла компьютерной сети}
\author{Приложение survivability}
\date{}

\begin{document}
\maketitle

\section{Постановка задачи}

Система описывается структурной моделью $S = \langle E, L, F \rangle$, где
$E = \{e_i\}$~--- множество элементов (узлов), $n$~--- их число,
$L$~--- связи, $F$~--- целевая функция системы.

На узел $e_i$ действуют деструктивные воздействия
$D = \{d_1, \dots, d_m\}$ (случайные отказы, атаки, разрывы каналов) с
интенсивностями $\lambda_j$. Каждое воздействие $d_j$ поражает узел $e_i$ с
вероятностью
\begin{equation}
p_{ij} = P\!\left(x_i = 0 \mid d_j\right),
\end{equation}
где $x_i \in \{0, 1\}$~--- состояние узла (1~--- работоспособен, 0~--- отказ).

Требуется за период оценивания $T$ найти \emph{живучесть узла}~--- вероятность
сохранения работоспособности к моменту $T$ при условии, что в начальный момент
узел работоспособен:
\begin{equation}
Q_i = P\!\left(x_i(T) = 1\right).
\end{equation}

\section{Модель деструктивных воздействий}

Собственные отказы узла и каждое деструктивное воздействие образуют
пуассоновские потоки событий. Воздействие $d_j$ реализуется как отказ узла с
интенсивностью $\lambda_j$ и каждая реализация независимо поражает узел с
вероятностью $p_{ij}$. Прореживание пуассоновского потока с вероятностью
$p_{ij}$ даёт пуассоновский поток поражений с интенсивностью $\lambda_j p_{ij}$.

Сумма независимых пуассоновских потоков~--- пуассоновский поток, поэтому
общая интенсивность ущерба узла:
\begin{equation}
\Lambda_i = \lambda_i^{(0)} + \sum_{j=1}^{m} \lambda_j \, p_{ij},
\label{eq:lambda}
\end{equation}
где $\lambda_i^{(0)}$~--- интенсивность собственных (случайных) отказов узла.

\section{Марковская модель «отказ--восстановление»}

Функционирование узла описывается однородной марковской цепью с
непрерывным временем и двумя состояниями:
\begin{itemize}
  \item состояние 1~--- узел работоспособен;
  \item состояние 0~--- узел отказал (поражён).
\end{itemize}

Переход $1 \to 0$ происходит с суммарной интенсивностью отказов и атак
$\Lambda_i$ (формула~\eqref{eq:lambda}), переход $0 \to 1$ (восстановление)~---
с интенсивностью $\mu \geq 0$ (экспоненциальное время восстановления;
$\mu = 0$ соответствует отсутствию восстановления).

Генератор марковской цепи:
\begin{equation}
A = \begin{pmatrix}
-\Lambda_i & \Lambda_i \\
\mu & -\mu
\end{pmatrix}.
\end{equation}

\subsection{Случай без восстановления ($\mu = 0$)}

Вероятность безотказной работы за время $T$:
\begin{equation}
Q_i(T) = e^{-\Lambda_i T}.
\label{eq:cold}
\end{equation}

\subsection{Случай с восстановлением ($\mu > 0$)}

Решение уравнения Колмогорова для вероятности работоспособности
$P_1(t) = P(x_i(t) = 1)$ при начальном условии $P_1(0) = 1$:
\begin{equation}
Q_i(T) = P_1(T) = \frac{\mu}{\Lambda_i + \mu}
+ \frac{\Lambda_i}{\Lambda_i + \mu}\, e^{-(\Lambda_i + \mu) T}.
\label{eq:main}
\end{equation}

Предельный (стационарный) коэффициент готовности:
\begin{equation}
K = \lim_{T \to \infty} Q_i(T) = \frac{\mu}{\Lambda_i + \mu}.
\label{eq:steady}
\end{equation}

Формула~\eqref{eq:main} имеет наглядный смысл: за время $T$ узел с
вероятностью $e^{-(\Lambda_i+\mu)T}$ ни разу не сменил состояние и остался
работоспособным (второе слагаемое), либо цепь успела достичь стационарного
режима и узел работоспособен с вероятностью $K$ (первое слагаемое).

Частные случаи:
\begin{align}
T \to 0 &: \quad Q_i(T) \to 1; \\
\mu = 0 &: \quad Q_i(T) = e^{-\Lambda_i T} \quad \text{(формула~\eqref{eq:cold})}; \\
\Lambda_i = 0 &: \quad Q_i(T) = 1 \quad \text{(ущерб невозможен)}.
\end{align}

\section{Вклады источников ущерба}

Для выявления доминирующей угрозы каждому источнику ущерба сопоставляется его
доля в общей интенсивности:
\begin{equation}
s_j = \frac{\lambda_j p_{ij}}{\Lambda_i}, \qquad
s_0 = \frac{\lambda_i^{(0)}}{\Lambda_i}, \qquad
s_0 + \sum_{j=1}^{m} s_j = 1.
\end{equation}

При $\Lambda_i = 0$ доли полагаются равными нулю.

\section{Оценка живучести системы}

Задача 2 слайда 8 презентации. Система описывается структурной моделью
$S = \langle E, L, F \rangle$; вектор состояний
$X = (x_1, \dots, x_n) \in \Omega = \{0, 1\}^n$ задаёт состояние всех узлов.
Состояния узлов принимаются независимыми, поэтому вероятность вектора
состояний равна
\begin{equation}
P(X) = \prod_{i=1}^{n} Q_i^{x_i} (1 - Q_i)^{1 - x_i},
\label{eq:px}
\end{equation}
где $Q_i$~--- живучесть узла $e_i$ за период $T$ по формуле~\eqref{eq:main}
(задача~1).

Функция эффективности $W(X) \in [0, 1]$ оценивает долю выполнения целевой
функции $F$ при состоянии $X$; система допустима, если
$W(X) \geq W_{\text{доп}}$, где $W_{\text{доп}}$~--- допустимый порог
деградации. \emph{Живучесть системы}~--- вероятность допустимого состояния:
\begin{equation}
Q(S) = P\!\left(W(X) \geq W_{\text{доп}}\right)
     = \sum_{X:\, W(X) \geq W_{\text{доп}}} P(X).
\label{eq:qs}
\end{equation}

Реализованы два вида целевой функции (поле \texttt{target} входного JSON):

1. доля работоспособных узлов (связи не учитываются):
\begin{equation}
W(X) = \frac{1}{n} \sum_{i=1}^{n} x_i;
\label{eq:wfraction}
\end{equation}
2. доля пар узлов, между которыми существует путь по работоспособным узлам
и связям (по компонентам связности работоспособной части графа):
\begin{equation}
W(X) = \frac{\sum_k \binom{s_k}{2}}{\binom{n}{2}},
\label{eq:wconn}
\end{equation}
где $s_k$~--- размер $k$-й компоненты связности; при $n = 1$ полагается
$W = x_1$.

\subsection{Точный перебор}

При $n \leq 20$ сумма~\eqref{eq:qs} вычисляется прямым перебором всех
$2^n$ векторов состояний. Вырожденные случаи обрабатываются естественно:
при всех $Q_i = 1$ остаётся только полностью работоспособный вектор и
$Q(S) = 1$; при всех $Q_i = 0$~--- $Q(S) = 0$ при $W_{\text{доп}} > 0$
(система не выполняет целевую функцию); при пустом множестве связей $L$
для функции~\eqref{eq:wconn} имеем $W(X) = 0$ при $n \geq 2$.

\subsection{Метод Монте-Карло}

При $n > 20$ выполняется $N$ независимых прогонов; в каждом состояние узла
$i$ генерируется по схеме Бернулли с вероятностью $Q_i$ (общий поток
случайных чисел с фиксированным зерном~--- см. раздел про задачу~3).
Оценка $\hat{Q}(S) = N^{*} / N$, где $N^{*}$~--- число прогонов с
допустимым состоянием. Точность оценки задаётся доверительным интервалом
Уилсона с доверительной вероятностью 95~\% ($z = 1{,}96$):
\begin{equation}
\hat{Q}(S) \in
\left[
\frac{\hat{p} + \frac{z^2}{2N} \;-\; z\sqrt{\frac{\hat{p}(1-\hat{p})}{N}
+ \frac{z^2}{4N^2}}}{1 + \frac{z^2}{N}}
\;;\;
\frac{\hat{p} + \frac{z^2}{2N} \;+\; z\sqrt{\frac{\hat{p}(1-\hat{p})}{N}
+ \frac{z^2}{4N^2}}}{1 + \frac{z^2}{N}}
\right], \qquad \hat{p} = \frac{N^{*}}{N}.
\label{eq:wilson}
\end{equation}
Интервал Уилсона корректен и при малых $N^{*}$, в отличие от интервала на
основе нормальной аппроксимации.

\subsection{Важность узлов}

Для выявления узлов, отказ которых критичен, используется критичность
Бирнбаума~--- условная вероятность того, что система допустима при
гарантированно работающем узле и не допустима при гарантированно
отказавшем:
\begin{equation}
I_i = Q(S \mid x_i = 1) - Q(S \mid x_i = 0),
\qquad
Q(S \mid x_i = 1) = \frac{1}{Q_i}
\sum_{\substack{X:\, x_i = 1 \\ W(X) \geq W_{\text{доп}}}} P(X),
\qquad
Q(S \mid x_i = 0) = \frac{1}{1 - Q_i}
\sum_{\substack{X:\, x_i = 0 \\ W(X) \geq W_{\text{доп}}}} P(X).
\label{eq:birnbaum}
\end{equation}
Дополнительно вычисляется потенциал повышения живучести
$\Delta_i^{+} = Q(S \mid x_i = 1) - Q(S)$: прирост $Q(S)$ от идеального
узла. При точном перебре суммы в~\eqref{eq:birnbaum} накапливаются за один
проход по $\Omega$; при Монте-Карло условные вероятности оцениваются
$2\,000$ условными прогонами на узел.

\section{Выбор рационального способа защиты}

Задача 3 (блок «Выбор способа защиты» слайда 8, ЧНЗ №4 слайда~9).
Каждый вариант защиты $g \in G$ описывается изменёнными параметрами узлов
(сниженные $\lambda_i^{(0)}$, повышенное $\mu$, пониженные вероятности
поражения $p_{ij}$) и/или добавленными резервными связями, а также
стоимостью $C(g) \geq 0$. Вариант $g$ порождает производную систему
$S_g$, для которой оценивается $Q_g(S)$ по формулам~\eqref{eq:qs} тем же
способом (перебор либо Монте-Карло). Рациональный способ защиты~---
допустимый по бюджету вариант с максимальной живучестью:
\begin{equation}
g^{*} = \arg\max_{g \in G:\; C(g) \leq C_{\text{доп}}} Q_g(S).
\label{eq:argmax}
\end{equation}
Если множество $\{g : C(g) \leq C_{\text{доп}}\}$ пусто, решения не
существует и сообщается перечень стоимостей вариантов; при равенстве
$Q_g(S)$ выбирается вариант с меньшей стоимостью, затем с меньшим именем.

Для сопоставимости оценок Монте-Карло разных вариантов используется
\emph{общий поток случайных чисел}: все оценки (базовая система и все
варианты) выполняются последовательно над одним генератором с фиксированным
зерном. Оценки $\hat{Q}_g(S)$ зависимы, и их различия отражают различия
вариантов, а не случайность потоков; сравнение при этом консервативно,
что допустимо при обосновании решения. Зерно фиксируется во входном файле
и перекрывается опцией \texttt{{-}-seed}.

Обоснование решения сопровождается производными показателями:
\begin{equation}
\Delta Q = Q_{g^{*}}(S) - Q_{\text{base}}(S), \qquad
\frac{\Delta Q}{Q_{\text{base}}}, \qquad
\frac{\Delta Q}{C(g^{*})},
\label{eq:deltaq}
\end{equation}
где $Q_{\text{base}}(S)$~--- живучесть системы без защиты, $\Delta Q /
C(g^{*})$~--- эффективность затрат (прирост живучести на единицу стоимости).

\section{Расшифровка обозначений}

\begin{center}
\begin{tabular}{|c|l|}
\hline
Обозначение & Смысл \\
\hline
$e_i$ & оцениваемый узел сети \\
$\lambda_i^{(0)}$ & интенсивность собственных (случайных) отказов узла $e_i$ \\
$d_j$ & $j$-е деструктивное воздействие \\
$\lambda_j$ & интенсивность воздействия $d_j$ \\
$p_{ij}$ & вероятность поражения узла $e_i$ воздействием $d_j$ \\
$\Lambda_i$ & общая интенсивность ущерба узла (формула~\eqref{eq:lambda}) \\
$\mu$ & интенсивность восстановления ($\mu = 0$~--- восстановления нет) \\
$T$ & период оценивания \\
$x_i$ & состояние узла: 1~--- работоспособен, 0~--- отказ \\
$Q_i(T)$ & живучесть узла: $P(x_i(T) = 1)$ \\
$K$ & предельный коэффициент готовности (формула~\eqref{eq:steady}) \\
$s_j$ & доля источника ущерба в общей интенсивности $\Lambda_i$ \\
\hline
$S = \langle E, L, F \rangle$ & структурная модель системы (узлы, связи, целевая функция) \\
$X$ & вектор состояний системы, $X \in \Omega = \{0, 1\}^n$ \\
$P(X)$ & вероятность вектора состояний (формула~\eqref{eq:px}) \\
$W(X)$ & функция эффективности: доля выполнения целевой функции \\
$W_{\text{доп}}$ & допустимый порог деградации \\
$Q(S)$ & живучесть системы: $P(W(X) \geq W_{\text{доп}})$ (формула~\eqref{eq:qs}) \\
$s_k$ & размер $k$-й компоненты связности работоспособной части графа \\
$N$, $N^{*}$ & число прогонов Монте-Карло (успешных) \\
$\hat{Q}(S)$ & оценка $Q(S)$ методом Монте-Карло, $N^{*}/N$ \\
$z$ & квантиль $N(0, 1)$: $1{,}96$ для доверительной вероятности 95~\% \\
$I_i$ & критичность Бирнбаума узла $e_i$ (формула~\eqref{eq:birnbaum}) \\
$G$ & множество вариантов способа защиты \\
$g$, $C(g)$ & вариант защиты и его стоимость \\
$C_{\text{доп}}$ & бюджет на защиту \\
$g^{*}$ & рациональный способ защиты (формула~\eqref{eq:argmax}) \\
$Q_g(S)$ & живучесть системы с защитой $g$ \\
$Q_{\text{base}}(S)$ & живучесть системы без защиты \\
$\Delta Q$ & прирост живучести от выбранного способа защиты (формула~\eqref{eq:deltaq}) \\
\hline
\end{tabular}
\end{center}

\section{Связь с формальной постановкой из презентации}

Задача 1 слайда 8 презентации («Оценка живучести узла») формулируется как:
дано $e_i$, $\lambda_i^{(0)}$, воздействия $d_j$ с интенсивностями $\lambda_j$,
вероятности поражения $p_{ij}$ и период $T$; найти
$Q_i = P(x_i(T) = 1)$. Формула~\eqref{eq:main} даёт явное решение этой задачи
в рамках марковской модели «отказ--восстановление» (направление «Теория
надёжности и резервирования» слайда~6). Совместная динамика отказов и атак
учитывается естественным образом через суммарную интенсивность
$\Lambda_i$~--- в отличие от раздельного учёта, она не завышает оценку
живучести.

Задача 2 («Оценка живучести системы $S$») решается формулами
\eqref{eq:px}--\eqref{eq:wilson}: при независимых узлах вероятность вектора
состояний факторизуется, а сумма по допустимым состояниям вычисляется
перебором ($n \leq 20$) или методом Монте-Карло с доверительным интервалом
Уилсона ($n > 20$). Это и есть оценка показателя живучести системы
прогонами на её цифровом двойнике (ЧНЗ №3): модель системы с теми же
параметрами прогоняется многократно, успешные прогоны дают $\hat{Q}(S)$.

Блок «Выбор способа защиты» слайда 8 ($g^{*} = \arg\max_{g \in G} Q_g(S)$
при $C(g) \leq C_{\text{доп}}$) реализован формулой~\eqref{eq:argmax}:
каждый вариант защиты оценивается прогонами на цифровом двойнике (ЧНЗ №4),
а сравнение вариантов выполняется на общем потоке случайных чисел, что
устраняет научное противоречие слайда~7 (оценки на единичных сценариях
невоспроизводимы): результат детерминирован при фиксированном зерне и
сопровождается доверительными интервалами.

\end{document}
